\section{\texttt{printk} 实现概述}
\label{sec:printk}

\texttt{printk} 是内核的 \texttt{printf}——唯一的格式化输出设施。
所有输出最终通过 \texttt{serial\_putc} 发送到串口。

\subsection{核心结构}

\codefile{src/kernel/core/printk.c}
\begin{lstlisting}[style=cstyle]
void printk(const char* fmt, ...) {
    va_list args;
    va_start(args, fmt);
    vprintk(fmt, args);     // 内部实现
    va_end(args);
}
\end{lstlisting}

\texttt{vprintk} 逐字符扫描 \texttt{fmt}，遇到 \texttt{\%} 时解析格式说明符。
支持的格式如表~\ref{tab:printk-formats}：

\begin{table}[H]
\centering
\begin{tabular}{llp{6cm}}
\toprule
\textbf{说明符} & \textbf{参数类型} & \textbf{说明} \\
\midrule
\texttt{\%d} & \texttt{int}      & 有符号十进制 \\
\texttt{\%u} & \texttt{unsigned} & 无符号十进制 \\
\texttt{\%x} & \texttt{unsigned} & 小写十六进制（无 \texttt{0x} 前缀） \\
\texttt{\%p} & \texttt{void*}    & 指针，格式 \texttt{0xXXXXXXXX}（固定 8 位） \\
\texttt{\%s} & \texttt{char*}    & 字符串（NULL 安全，输出 \texttt{(null)}） \\
\texttt{\%c} & \texttt{int}      & 单个字符 \\
\texttt{\%\%} & —               & 字面 \texttt{\%} \\
\bottomrule
\end{tabular}
\caption{\texttt{printk} 支持的格式说明符}
\label{tab:printk-formats}
\end{table}

\subsection{格式化选项}

\texttt{vprintk} 还支持：
\begin{itemize}
  \item \textbf{宽度} —— \texttt{\%8x} 输出至少 8 个字符（右对齐，空格填充）。
  \item \textbf{零填充} —— \texttt{\%08x} 左侧补零。
  \item \textbf{左对齐} —— \texttt{\%-16s} 左对齐、右侧补空格。
\end{itemize}

\begin{designbox}
为什么从头实现 \texttt{printk} 而不用 \texttt{libc} 的 \texttt{printf}？

内核没有标准 C 库（freestanding 环境）。\texttt{printf} 依赖 \texttt{FILE*}、
\texttt{stdout}、\texttt{write} 系统调用等基础设施——这些都是我们自己要实现的东西。
内核的 \texttt{printk} 直接写 \texttt{serial\_putc}，零依赖、零缓冲。
\end{designbox}

% ============================================================================

